\section{Preliminaries}
\label{sec:prelim}

\subsection{Logic Synthesis and the EDA Flow}
\label{sec:prelim:synthesis}

\todo{lets call this subsection Electronic Design Automation and Logic Synthesis}
\todo{add subsubseciton on EDA itself what is it what is it used for why should i care}

\subsubsection{Position of Logic Synthesis in the Design Flow}
\label{sec:prelim:synthesis:flow}
Electronic Design Automation (EDA) organises the design of a digital circuit as a flow of
successive refinements \citep{demicheli1994synthesis}. A Register-Transfer-Level (RTL)
description is compiled into a network of Boolean gates, Logic Synthesis (LS) restructures
that network against cost objectives, technology mapping binds the result to the cells of
a target library, and placement and routing assign physical positions and wires. LS is
pre-physical: the network it operates on has no coordinates, no wire lengths and no
layout-extracted timing, which is what makes the spatial coarsening methods
of~\ref{sec:relwork:domain:circuits} inapplicable to it.


\todo{consider shoudl these subsubseciotn not be after mayeb introducing AIGs and with the algorithm.tex }
\subsubsection{Optimization Objectives: Area, Power, Delay}
\label{sec:prelim:synthesis:objectives}
Synthesis quality is measured against three costs: the silicon area a circuit occupies,
the power it draws, and the delay of its longest combinational path. None of the three is
physical before technology mapping, so synthesis tracks each through a structural proxy,
gate count for area and logic depth for delay, both formalized on the AIG representation
in~\ref{sec:prelim:aig}. This work uses node count as its area proxy and models nothing
else (\ref{sec:intro:scope:target}). The validity cost of that restriction is assessed
in~\ref{sec:discussion:limitations:validity}. \todo{mentioned already earlier,prelim should just be introduction to concepts not arguments abotu psitions and thesis}

\subsubsection{Synthesis Scripts and Pipelines}
\label{sec:prelim:synthesis:scripts}
A synthesis script is an ordered sequence of function-preserving transformations applied
as a composition,
\begin{equation}
    S = t_m \circ t_{m-1} \circ \cdots \circ t_1,
    \label{eq:prelim:script}
\end{equation}
and order matters: the $t_i$ do not commute, and the gain of each depends on the structure
its predecessors produced, so the same script yields different gains on different designs.
Script choice can therefore not be fixed once and reused across designs. The scripts used
in this work are the subject of~\ref{sec:prelim:algorithms}.
\todo{work in somehow the work recipe since it is used and steps and what they are as this is used to explain teh dataset generation}
\todo{also i dont htink this mentions why one needs and does AIG optimization like why smaller why is that better}